\BourbakiExerciseGroup

\begin{enumerate}
\def\labelenumi{\arabic{enumi})}
\tightlist
\item
  Let \((x_{t})\) be a family of real numbers, all of which belong to the interval \(A' = ] - \infty, +\infty]\) (resp. \(A'' = [-\infty, +\infty[\) ). The family \((x_{t})\) is summable in \(\overline{\mathbf{R}}\) if and only if one of the following conditions is satisfied:
\end{enumerate}

\begin{enumerate}
\def\labelenumi{\alph{enumi})}
\item
  At least one of the numbers \(x_{t}\) is equal to \(+\infty\) (resp. \(-\infty\) ).
\item
  At least one of the two sums \(\sum_{i} x_{i}^{+}, \sum_{i} x_{i}^{-}\) is finite.
\end{enumerate}

In the first case we have \(\sum_{i} x_{i} = +\infty\) (resp. \(-\infty\) ); in the second case, \(\sum_{i} x_{i} = \sum_{i} x_{i}^{+} - \sum_{i} x_{i}^{-}\) .

\begin{enumerate}
\def\labelenumi{\arabic{enumi})}
\setcounter{enumi}{1}
\item
  Let \((x_{i})\) be a family of real numbers all of which belong to the interval \(\mathbf{A}'\) (resp. \(\mathbf{A}''\) ) and which satisfies condition \(b\) ) of Exercise 1. Show that every subfamily of \((x_{i})\) is summable in \(\overline{\mathbf{R}}\) , and that the sum of the \(x'\) is associative (cf.~Chapter III, § 5, no. 3, Theorem 2).
\item
  Let \((x_{n})_{n\geqslant 0}\) be a decreasing sequence of finite real numbers \(\geqslant 0\) . This sequence is summable in \(\mathbf{R}\) if and only if the sequence \((2^{n}x_{2^{n}})_{n\geqslant 0}\) is summable in \(\mathbf{R}\) (``Cauchy's condensation test''). * Deduce that the series whose general term is \(1 / n^{\alpha}\) is convergent if \(\alpha >1\) and not convergent if \(0 < \alpha \leqslant 1\) . *
\item
  Let \((a_{n})\) be any sequence of finite real numbers \(\geqslant 0\) . Show that the sequence \(\left(\frac{a_n}{(1 + a_1)(1 + a_2)\ldots(1 + a_n)}\right)\) is summable in \(\mathbf{R}\) (express each term of this sequence as a difference).
\item
  Let \((d_n)\) be a sequence of finite real numbers \(\geqslant 0\) such that
\end{enumerate}

\[
\sum_ {n = 0} ^ {\infty} d _ {n} = + \infty .
\]

What can be said about the convergence of the series whose general terms are

\[
\frac {d _ {n}}{\mathrm{I} + d _ {n}}, \quad \frac {d _ {n}}{\mathrm{I} + n d _ {n}}, \quad \frac {d _ {n}}{\mathrm{I} + n ^ {2} d _ {n}}, \quad \frac {d _ {n}}{\mathrm{I} + d _ {n}}?
\]

\begin{enumerate}
\def\labelenumi{\arabic{enumi})}
\setcounter{enumi}{5}
\item
  Prove that \(s = \sum_{n=1}^{\infty} (1/n)\) is equal to \(+\infty\) by showing that \(s \geqslant s + \frac{1}{2}\) (find lower bounds, as functions of \(s\) , for the sum of the terms with even indices and the sum of the terms with odd indices).
\item
  Show that, for each integer \(p > \mathbf{r}\) ,
\end{enumerate}

\[
\sum_ {n = 1} ^ {\infty} \frac {(- \mathrm{I}) ^ {n - 1}}{n ^ {p}} = (\mathrm{I} - 2 ^ {1 - p}) \sum_ {n = 1} ^ {\infty} \frac {\mathrm{I}}{n ^ {p}}.
\]

\begin{enumerate}
\def\labelenumi{\arabic{enumi})}
\setcounter{enumi}{7}
\item
  Show that, if \(m\) is an integer \(> 0\) , then \(\sum_{n \geqslant 1, n \neq m} \frac{1}{m^2 - n^2} = -\frac{3}{4m^2}\) (express the rational fraction \(\frac{1}{m^2 - x^2}\) as a sum of partial fractions).
\item
  If the series whose general term is \(x_{n}\) is convergent in \(\mathbf{R}\) , show that \(\liminf_{n\to \infty}nx_n\leqslant 0\leqslant \limsup_{n\to \infty}nx_n\) .
\end{enumerate}

¶ 10) A series \((u_n)\) is convergent in \(\mathbb{R}\) if and only if, for each increasing sequence \((p_n)\) of numbers \(>0\) such that \(\lim_{n\to \infty}p_n = +\infty\) , we have \(\lim_{n\to \infty}\left(\left(\sum_{k=0}^{n}p_ku_k\right) / p_n\right) = 0\) (use Exercise 15 of § 5).

\begin{enumerate}
\def\labelenumi{\arabic{enumi})}
\setcounter{enumi}{10}
\tightlist
\item
  Consider a sequence \((a_{n})\) every term of which is the sum of a finite sequence of finite real numbers
\end{enumerate}

\[
a _ {n} = b _ {n, 1} + b _ {n, 2} + \dots + b _ {n, k _ {n}}.
\]

For each pair \((n, p)\) of positive integers such that \(p \leqslant k_n\) , if \(m = \sum_{i=0}^{\infty} k_i + p\) , put \(c_m = b_{n,p}\) . Show that if the series whose general term is \(a_n\) is convergent in \(\mathbf{R}\) , and if

\[
d _ {n} = \left| b _ {n, 1} \right| + \left| b _ {n, 2} \right| + \dots + \left| b _ {n, k _ {n}} \right|
\]

tends to o as n tends to infinity, then the series whose general term is \(c_{m}\) is convergent and has the same sum as the series whose general term is \(a_{n}\) .

¶ 12) Let \((u_n)\) be a sequence of finite real numbers. For each permutation \(\sigma\) of the set N, put

\[
r (n) = | \sigma (n) - n |. \sup _ {m \geqslant n} | u _ {m} |.
\]

\begin{enumerate}
\def\labelenumi{\alph{enumi})}
\item
  If the series with general term \(u_n\) is convergent, and if \(\lim_{n \to \infty} r(n) = 0\) , show that the series with general term \(v_n = u_{\sigma(n)}\) is convergent to the same sum. {[}Consider the difference \(\sum_{k=0}^{n} v_k - \sum_{k=0}^{n} u_k\) for large values of \(n\) , and if \(h\) is the smallest integer \(\leqslant n\) such that \(\sigma(h) > n\) , note that in this difference there are at most \(\sigma(h) - h\) terms in each of the two sums which do not cancel out{]}.
\item
  Suppose that the series whose general term is \(u_n\) is convergent. Give an example of a permutation \(\sigma\) such that \(\lim_{n \to \infty} r(n) = 0\) but such that the criterion of Chapter III, § 5, Exercise 6 a) is not satisfied. {[}Define a family of mutually disjoint intervals \(I_k = [n_k - 2k, n_k + 2k]\) in \(N\) , and take \(\sigma\) such that \(\sigma(n) = n\) whenever \(n\) does not belong to any of the \(I_k\) , and such that \(\sigma(n_k - 2j) = n_k + 2j\) , \(\sigma(n_k + 2j) = n_k - 2j\) for all \(k\) and \(0 \leqslant j \leqslant k\) .{]}
\end{enumerate}

\begin{enumerate}
\def\labelenumi{\arabic{enumi})}
\setcounter{enumi}{12}
\item
  Let \((u_n)\) be a sequence of real numbers \(\geqslant 0\) such that \(\lim_{n\to \infty}u_n = 0\) and \(\sum_{n = 0}^{\infty}u_n = +\infty\) . Let \((p_n)\) be a strictly increasing sequence of numbers \(>0\) such that \(\lim_{n\to \infty}p_n = +\infty\) . Show that there exists a permutation \(\sigma\) of \(\mathbf{N}\) such that for each \(n\in \mathbf{N}\) , \(\sum_{k = 0}^{n}u_{\sigma (n)}\leqslant p_n\)
\item
  Let \((u_{n})\) be a sequence of finite real numbers such that the series whose general term is \(u_{n}\) is convergent but not absolutely convergent; let \(s = \sum_{n=0}^{\infty} u_{n}\) . Show that, for each number \(s' \geqslant s\) , there exists a permutation \(\sigma\) of \(\mathbf{N}\) such that \(\sigma(n) = n\) for those indices \(n\) such that \(u_{n} \geqslant 0\) , and such that \(\sum_{n=0}^{\infty} u_{\sigma(n)} = s'\) . {[}Show by induction on \(m\) that there is a permutation \(\sigma_m\) of \(\mathbf{N}\) such that \(\sigma_m(k) = k\) for all \(k\) for which \(u_k \geqslant 0\) and such that, if \(u_k^{(m)}\) denotes \(u_{\sigma_m(k)}\) , there is an integer \(p_m\) with the property that \(\left| s' - \sum_{i=0}^{k} u_i^{(m)} \right| \leqslant 1/m\) for all \(k \geqslant p_m\) . Moreover, \(\sigma_{m+1}\) is such that \(\sigma_{m+1}(k) = \sigma_m(k)\) for all \(k\) such that \(\sigma_m(k) < p_m\) and all \(k\) such that \(u_k \leqslant -1/m\) .{]}
\item
  Let \((u_{n})\) be a sequence of finite real numbers such that \(\lim_{n\to \infty}u_n = 0\) and \(\sum_{n}u_{n}^{+} = \sum_{n}u_{n}^{-} = +\infty\) . If \(a\) and \(b\) are any two real numbers (finite or not) such that \(a\leqslant b\) , show that there is a permutation \(\sigma\) of \(\mathbf{N}\) such that, if we put \(s_n = \sum_{k = 0}^n u_{\sigma (k)}\) , we have \(\liminf_{n\to \infty}s_n = a\) and \(\limsup_{n\to \infty}s_n = b\) . Show that the set of cluster points of the sequence \((s_n)\) is then the interval \([a,b]\) .
\item
  Let \((u_n)\) be a sequence of finite real numbers such that the series whose general term is \(u_n\) is convergent but not absolutely convergent. Show that there exists a permutation \(\sigma\) of \(\mathbf{N}\) satisfying the condition of Chapter III, § 5, Exercise 6 \(a\) ) but such that the series whose general term is \(u_{\sigma - 1(n)}\) is not convergent. {[}Let \(h, k, m\) be three integers with the following properties: if \(p_1, \ldots, p_r\) are the integers \(n\) such that \(h \leqslant n < h + m\) and \(u_n \geqslant 0\) , then \(u_{p_1} + \cdots + u_{p_r} > 2\) and \(h + m < k - r\) ; moreover if \(\mu, \nu\) are any integers such that \(h + m \leqslant \mu \leqslant \nu\) , then \(\left|\sum_{n = \mu}^{\nu} u_n\right| \leqslant 1\) . Put \(s = m - r\) and let \(q_1, \ldots, q_s\) denote the integers \(n\) such that \(h \leqslant n < h + m\) and \(u_n < 0\) . Consider the permutation \(\pi\) of the interval \([h, k + s]\) of \(\mathbf{N}\) such that \(\pi(p_i) = k - i + 1\) for \(i \leqslant i \leqslant r\) , \(\pi(q_j) = k + j\) for \(i \leqslant j \leqslant s\) , \(\pi(k + j) = h + s - j\) for \(i \leqslant j \leqslant s\) , \(\pi(k - i + 1) = h + s + i - 1\) for \(i \leqslant i \leqslant r\) and finally \(\pi(n) = n\) for \(h + m \leqslant n \leqslant k - r\) ; show that we have
\end{enumerate}

\[
\left| \sum_ {i = h} ^ {k} u _ {i} - \sum_ {i = h} ^ {k} u _ {\pi^ {- 1} (i)} \right| \geqslant I.
\]

\begin{enumerate}
\def\labelenumi{\arabic{enumi})}
\setcounter{enumi}{16}
\tightlist
\item
  If \(u_{mn} = \mathbf{i} / (m^2 - n^2)\) for \(m \neq n\) and \(u_{mm} = 0\) , show that
\end{enumerate}

\[
\sum_ {m = 1} ^ {\infty} \left(\sum_ {n = 1} ^ {\infty} u _ {m n}\right) = - \sum_ {n = 1} ^ {\infty} \left(\sum_ {m = 1} ^ {\infty} u _ {m n}\right) \neq 0.
\]

(Use Exercise 8.)

\begin{enumerate}
\def\labelenumi{\arabic{enumi})}
\setcounter{enumi}{17}
\item
  Let \((\mathbf{i} + u_{\mathfrak{i}})\) be a family of real numbers all belonging to the interval \([0, +\infty[\) (resp. \(]0, +\infty]\) ). Then the family \((\mathbf{i} + u_{\mathfrak{i}})\) is multipliable in \(\mathbb{R}\) if and only if one of the following conditions is satisfied: a) At least one of the numbers \(\mathbf{i} + u_{\mathfrak{i}}\) is equal to o (resp. \(+\infty\) ). b) The family \(u_{\mathfrak{i}}\) is summable in \(\overline{\mathbb{R}}\) .
\item
  Let \((\mathbf{I} + u_{t})_{t\in \mathbf{I}}\) be a multipliable family in \(\mathbb{R}\) . For each non-empty subset \(\mathbf{H}\) of \(\mathbf{I}\) , let \(v_{\mathbf{H}}\) denote \(\prod_{i\in \mathbf{H}}u_i\) . Show that the family \((v_{\mathbf{H}})_{\mathbf{H}\in \mathfrak{F}(\mathbf{I})}\) is summable in \(\mathbb{R}\) and that
\end{enumerate}

\[
\mathrm{I} + \sum_ {\mathbf {H}} v _ {\mathbf {H}} = \prod_ {i \in \mathbf {I}} (\mathrm{I} + u _ {i}).
\]

Is the converse true?

Deduce that, if \(-\mathbf{I} < x < \mathbf{I}\) , then \(\prod_{n=0}^{\infty} (\mathbf{I} + x^{2^n}) = \mathbf{I} / (\mathbf{I} - x)\) .

\begin{enumerate}
\def\labelenumi{\arabic{enumi})}
\setcounter{enumi}{19}
\item
  Let \((u_n)\) be a sequence of finite real numbers \(\geqslant 0\) such that \(u_0 > 0\) , and let \(s_n = \sum_{k=0}^{n} u_k\) for each \(n \geqslant 0\) . Then the sequence \((u_n)\) is summable in \(\mathbf{R}\) if and only if the sequence \((u_n / s_n)\) is summable in \(\mathbf{R}\) {[}apply Theorem 4 to the sequence \((1 - u_n / s_n)]\) . The same holds if the sequence \((u_n / s_n)\) is replaced by the sequence \((u_n / s_{n-1})\) .
\item
  Let \(u_{n} = (-\mathrm{I})^{n} / \sqrt{n}\) for \(n \geqslant 2\) . The product whose general factor is \(\mathrm{I} + u_{n}\) is not convergent, but the series whose general term is \(u_{n}\) is convergent.
\item
  For \(n \geqslant 2\) , let \(u_{2n-1} = -\mathrm{I}/\sqrt{n}\) and \(u_{2n} = (\mathrm{I} + \sqrt{n})/n\) . The product whose general factor is \(\mathrm{I} + u_n\) is convergent, but the series whose general term is \(u_n\) is not convergent.
\end{enumerate}

